\section*{Original Section 4.2.1: Construction of the Garland (pp.17:11--12)}
To fulfill our plan of the proof, described at the beginning of the section, we start constructing the garland by obtaining an exponentially big set of paths with the corresponding assignments. Let us remind that $|M|\le\nu_k\Delta r$ and $\rho$ is $\beta'$-mlce on $M$.
We say that triple $(p,U_p,\sigma_p)$ is $\gamma$-good iff:
\begin{itemize}\item $p$ is a path from the root of the branching program;
\item $U$ is a subset of edges such that corresponding variables are queried on $p$, so-called ``branching variables'';
\item $\sigma_p$ is a partial assignment such that:
\begin{itemize}\item $\sigma_p$ extends $\rho\cup\tau_p$;
\item $\sigma_p$ is a $\gamma$-mlce on $M\cup U_p$;
\item $\sigma_p$ $0.8$-satisfies $V_{\sigma_p}\setminus V_\rho$,
\end{itemize}
where $\tau_p$ is an assignment that corresponds to $p$.\end{itemize}
We maintain the set of $\beta'$-good triples $\mathcal P$ and an auxiliary set $\mathcal S$ of triples that appear in the set $\mathcal P$ at some moment during the process. In the beginning of our construction $\mathcal P:=\{(\emptyset,\emptyset,\rho)\}$ and $\mathcal S:=\mathcal P$.
We repeat the following process while we have at least one triple $(p,U_p,\sigma_p)\in\mathcal P$ such that $|U_p|\le\nu_k\Delta r$.
Consider the triple described above. Let $v$ be the end of $p$ and $x_e$ be the variable asked in $v$.
\begin{enumerate}
\item If $x_e$ was queried on $p$ we stop the process. In this case we return ``Repetition'' and we remember the path $p$.
\item Erase the triple $(p,U_p,\sigma_p)$ from $\mathcal P$.
\item If $\sigma_p(x_e)\in\{0,1\}$, then we continue along the edge $x_e=\sigma_p(x_e)$. Consider a path $p'$ that is the extension of $p$ along this edge, $U_{p'}:=U_p$ and $\sigma_{p'}:=\sigma_p$. Put $(p',U_{p'},\sigma_{p'})$ into $\mathcal P$ and $\mathcal S$ and repeat the process from the beginning.
\item If $\sigma_p(x_e)=*$, then it is a ``branching node'', and we call this step a branching step.
\begin{enumerate}[label=\alph*.]
\item Let $p'$ be a path obtained by continuing $p$ along the edge $x_e=0$, and $p''$ be a path obtained by continuing $p$ along the edge $x_e=1$.
\item $U_{p'}:=U_p\cup e$, $U_{p''}:=U_p\cup e$.
\item $\tau':=\sigma_p\cup\{x_e=0,x_{e'}=0\}$, $\tau'':=\sigma_p\cup\{x_e=1,x_{e'}=0\}$, where $x_{e'}$ is a dual edge.
\item $(p',U_{p'},\tau')$ is $(\beta'-0.01)$-good triple. We extend an assignment $\tau'$ to make this triple $\beta'$-good. For the formal statement see Lemma 18. Here we describe an idea. Let $R\subseteq E$ be a set of edges that are unassigned by $\tau'$ (or $\tau''$), and $B\subseteq V\setminus V_{\sigma_p}$ be the maximal set of vertices that satisfies:
\[|B|\le r;\qquad |E(B,\overline B)\cap R|\le\beta'\Delta|B|.\]
Let $\kappa$ be an assignment on variables that correspond to edges in the set $E(B)\setminus\supp(\tau')$ such that $\tau'\cup\kappa$ $0.8$-satisfies the constraints for all $v\in B$. This assignment $\kappa$ always exists (and moreover it is independent of the value of $x_e$, but we do not use this fact).
\item We denote $\sigma_{p'}:=\tau'\cup\kappa$, $\sigma_{p''}:=\tau''\cup\kappa$ and put $(p',U_{p'},\sigma_{p'})$ and $(p'',U_{p''},\sigma_{p''})$ into $\mathcal P$ and into $\mathcal S$.
\end{enumerate}\end{enumerate}
To conclude the construction we want to show the following claims.
\begin{description}
\item[Repetition case.] In the first case of the proof (if we have a repetition) we can reduce the problem to a lower bound on $(1,+(k-1))$-BP.
\item[Correctness.] The branching step can be done and triples $(p',U_{p'},\sigma_{p'})$ and $(p'',U_{p''},\sigma_{p''})$ satisfy the required properties.
\item[Garland extraction.] Among these paths we can find an $k$-garland $(a,b)$ and a locally consistent extension of $\rho$.
\end{description}
\subsection*{4.2.1.1 Correctness}
We show that if we have a triple $(p,U_p,\tau_p)$ which is $\beta'$-good then after processing it with our algorithm we also put in our sets $\beta'$-good triples. Let us formulate the general Lemma that helps us with it.
\begin{result}{Lemma 18}
Let $(p,U_p,\sigma_p)$ and $(q,U_q,\sigma_q)$ be $0.9$-good triples. Then there is an assignment $\kappa$ such that:
\begin{itemize}\item for any $\gamma$ that is an instance of $\sigma_p\joint\sigma_q$ an assignment $\gamma\cup\kappa$ is a $\beta'$-mlce on $\supp(\sigma_p)\cup\supp(\sigma_q)$;
\item $|\supp(\gamma\cup\kappa)|\le\nu_{k-1}\Delta r$.\end{itemize}
Moreover if $p=q$ then triple $(p,U_p,\sigma_p\cup\kappa)$ is $\beta'$-good.
\end{result}
\begin{proof}
The proof was motivated by the closure technique developed in [1,3]. For the full version of the proof see Appendix A.
\end{proof}
If the branching step was not done, then we do not change $U$ and $\tau$, and we extend the path $p$ according to the assignment $\tau$ hence the triple remains $\beta'$-good. We are left with the branching step. Note that $(p',U_{p'},\tau')$ is $0.9$-good and we apply Lemma 18 to a pair composed of two identical triples $(p',U_{p'},\tau')$ and obtain $\kappa$ that satisfies the required properties.
\subsection*{4.2.1.2 Repetition case}
First let us note that if there is a repetition, then $k>0$. Suppose we found a repetition while considering a triple $(p,U_p,\sigma_p)$. The size of $M\cup U_p$ is at most $2\nu_k\Delta r$ and $\sigma_p$ is $\beta'$-mlce on $M\cup U_p$. Let $v$ be an end node of $p$. The program $D(v)|_{\sigma_p}$ is a $(1,+(k-1))$-BP for $\Search_{\Flow_G^{\eta^k}|_{\sigma_p}}$ by Lemma 8. Thus by induction hypothesis we have a lower bound of $2^{\frac{\nu_{k-1}}{4k^2}\Delta r}\ge2^{\frac{\nu_k}{4(k+1)^2}\Delta r}$ on the size of $D(v)|_{\sigma_p}$ and in this case we are done.
\subsection*{4.2.1.3 Garland extraction}
The following Lemma gives us a pair of triples $(p,U_p,\sigma_p),(q,U_p,\sigma_q)\in\mathcal P$ such that $(p,q)$ forms a $(k+1)$-garland.
\begin{result}{Lemma 19}
There are $(p,U_p,\sigma_p),(q,U_q,\sigma_q)\in\mathcal S$ such that $(p,q)$ forms a $(k+1)$-garland.
\end{result}
\begin{proof}For the proof see Appendix B.\end{proof}
To continue the proof we need some additional property that we can ``avoid repetitions'' in this garland. We say that there is a repetition in a garland $p=v_0p_1v_1p_2v_2p_3\ldots p_{k_0+1}v_{k_0+1}$ and $q=v_0q_1v_1q_2v_2q_3\ldots q_{k_0+1}v_{k_0+1}$ iff there is path in the garland, i.e. path $r$ of the form $v_0r_1v_1r_2v_2r_3\ldots r_{k_0+1}v_{k_0+1}$, such that some variable is queried more than one time on it, where $r_i\in\{p_i,q_i\}$.
Consider a path $r$ in our garland $(p,q)$ that contains a repetition and $r'\subseteq r$ the largest initial segment of $r$ without repetitions. Let $v$ be its end node. We apply Lemma 18 to triples $(p,U_p,\sigma_p),(q,U_q,\sigma_q)$, which gives us assignment $\kappa$, and choose a instance $\gamma$ of $\sigma_p\joint\sigma_q$ that is consistent with $\tau_{r'}$. Moreover, $|\supp(\gamma\cup\kappa)|\le\nu_{k-1}\Delta r$, and $\gamma\cup\kappa$ is a $\beta'$-mlce on $\supp(\sigma_p)\cup\supp(\sigma_q)$. Hence by Lemma 8 we can use the induction hypothesis for $(1,+k-1)$-BP $D(v)|_{\tau_{r'}}$ and formula $\varphi|_{\gamma\cup\kappa}$. The size of $D(v)|_{\tau_{r'}}$ is at least $2^{\frac{\nu_{k-1}}{4k^2}\Delta r}\ge2^{\frac{\nu_k}{4(k+1)^2}\Delta r}$.
For the rest of the proof we can assume that on any path $r$ of the form described above there are no repetitions.
\section*{Original Section 4.2.2: Unreachable Leaves (pp.17:13--14)}
Let us summarize what we have from the previous section. We created a pair of triples: $(p,U_p,\sigma_p)$ and $(q,U_q,\sigma_q)$ such that:
\begin{itemize}\item $(p,q)$ forms $(k_0+1)$-garland:
\[p=v_0p_1v_1p_2v_2p_3\ldots p_{k_0+1}v_{k_0+1};\quad q=v_0q_1v_1q_2v_2q_3\ldots q_{k_0+1}v_{k_0+1};\]
\item $(p,U_p,\sigma_p)$ and $(q,U_q,\sigma_q)$ are $\beta'$-good;
\item there are no repetitions on any path in the garland $(p,q)$.\end{itemize}
We use Lemma 18 for $(p,U_p,\sigma_p)$ and $(q,U_q,\sigma_q)$ and get an assignment $\kappa$. Let us fix an assignment $\gamma$ that is an instance of $\sigma_p\joint\sigma_q$ consistent with:
\begin{itemize}\item $\tau_p$;
\item values in $\sigma_q$ that do not contradict $\tau_p$\end{itemize}
and denote $\zeta:=\gamma\cup\kappa$. Note that, by construction:
\begin{itemize}\item $|\zeta|\le\nu_{k-1}\Delta r$;
\item $|\zeta|$ is $(r,0.6,\beta')$-locally consistent.\end{itemize}
In this section we describe a set of clauses that should be unreachable from the vertex $v_{k_0+1}$. Note that on each segment of a garland $(v_ip_iv_{i+1},v_iq_iv_{i+1})$ we query at least one variable in both assignments $\tau_p$ and $\tau_q$ and get the different values. Denote any variable that satisfies this property by $x_i$.
We remind that $\varphi:=\Flow_G^{\eta^{k_0+1}}$. Let $\mathcal D,\mathcal C$ be the subsets of clauses:
\[\mathcal D:=\{D\in\Flow_G\mid\text{for every }e\text{ that corresponds to some }x_i:e\in\eta^{k_0+1}(D)\}.\]
and
\[\mathcal C:=\{C\in\varphi\mid C\text{ is obtained from some }D\in\mathcal D\text{ by the amplification trick}\}.\]
For the sake of contradiction suppose that there is a path $s$ from $v_{k_0+1}$ such that:
\begin{itemize}\item $s$ is a consistent path and $\tau_s$ is consistent with $\zeta$ and hence $ps$ is also consistent;
\item $s$ ends in a clause $C\in\mathcal C$.\end{itemize}
Consider a family of paths $r_i:=v_0p_1v_1p_2v_2p_3\ldots p_{i-1}v_{i-1}q_iv_ip_{i+1}q_{i+1}p_{i+2}\ldots p_{k_0+1}v_{k_0+1}$, where $i\in[k_0+1]$. All paths $r_i$ are consistent since there are no repetitions in the garland $(p,q)$. Hence if $r_i$ is inconsistent with $s$ then on $s$ we requery some variable $x'_i$ from the segment $v_{i-1}q_iv_i$ and get an inconsistent value.
By construction, $\tau_s$ is consistent with $\zeta$, and $\zeta:=\gamma\cup\kappa$, where $\gamma$ is an instance of $\sigma_p\joint\sigma_q$. If $x'_i$ appeared in $v_{i-1}q_iv_i$, but not in $v_{i-1}p_iv_i$ (note that it cannot appear in any other segment of the garland, since there are no repetitions on the garland), then $(\sigma_p\joint\sigma_q)(x'_i)\in\{\sigma_q(x'_i),?\}$ and $\tau_p(x'_i)=*$ thus $\gamma(x'_i)=\sigma_q(x'_i)$ by the choice of $\gamma$. It follows that $\zeta(x'_i)=\sigma_q(x'_i)$ as well, and since $\tau_s$ is consistent with $\zeta$, we cannot obtain an inconsistent with $\tau_{q_i}$ value for $x'_i$ while requerying it. Hence $x'_i$ had appeared in $v_{i-1}p_iv_i$ as well, and on $s$ we requeried a variable from $v_{i-1}p_iv_i$ in consistent way. Moreover if all paths from some set $\{r_i\}_{i\in L}$ where $L\subseteq[k_0+1]$ are inconsistent with $s$ we requery at least $|L|$ variables from the path $p$ on the path $s$. Hence at least one of the paths $r_{i_0}$ is consistent with $s$, where $i_0\in[k_0+1]$ (or on the path $ps$ we requery at least $k_0+1$ variables).
\begin{result}{Remark 20}This is the only place there we use the property that there are no repetitions on the garland.\end{result}
Consider two paths $ps$ and $r_{i_0}s$:
\begin{itemize}\item these paths are consistent;
\item $\tau_{ps}(x_{i_0})\ne\tau_{r_{i_0}s}(x_{i_0})$.\end{itemize}
These properties imply that clause $C$ is not a legal answer for at least one these paths, and we have a contradiction with the assumption that there is a consistent path from $v_{k_0+1}$ to this clause. That gives us the desired description of leaves that should be unreachable for $v_{k_0+1}$.
To conclude the proof it remains to show that there should be a path from $v_{k_0+1}$ to at least one leaf marked by a clause $C\in\mathcal C$. We do it in the next section.
